\section*{Capítulo \ref{ch:b2:wave-interfaces} --- Reflexão e transmissão em interfaces}

\begin{solution}{exo:b2:wave-interfaces:1}
$r = (Z_2 - Z_1)/(Z_2 + Z_1)$, $R = r^2$: ar--água $0.9995$, $99.9\%$;
água--aço $0.93$, $86\%$; tecido--osso $0.63$, $40\%$; tecido--ar
$-0.9995$, $99.9\%$. Sua voz, no ar, é refletida pela superfície; o
motor sacode o casco, que é de aço em contato com a água e
passa $14\%$ de sua vibração direto para dentro.
\end{solution}

\begin{solution}{exo:b2:wave-interfaces:2}
$R = ((n - 1)/(n + 1))^2$: $2.0\%$, $4.3\%$, $17\%$; vidro--água $0.44\%$.
$0.957^{10} = 0.64$: um terço da luz perdido, boa parte como fantasmas;
com revestimento: $0.995^{10} = 0.95$.
\end{solution}

\begin{solution}{exo:b2:wave-interfaces:3}
$\tan\theta_B = n_2/n_1$: \qty{56.7}{\degree}, \qty{53.1}{\degree}, \qty{48.8}{\degree}. A
luz refletida fica polarizada perpendicularmente ao plano de incidência;
a luz transmitida fica apenas parcialmente polarizada (a componente s é
só parcialmente refletida). Olhe para a vitrine a \qty{57}{\degree} de sua
normal, com o eixo do polarizador no plano de incidência (horizontal,
para uma vitrine vertical vista de lado).
\end{solution}

\begin{solution}{exo:b2:wave-interfaces:4}
$\arcsin(1/n)$: \qty{41.1}{\degree}, \qty{48.8}{\degree}, \qty{24.4}{\degree}; $45^\circ > 41^\circ$:
o prisma reflete totalmente. $\ell = \lambda/2\pi\sqrt{n^2\sin^2\theta - 1}$: \qty{270}{nm}
a \qty{45}{\degree}, \qty{115}{nm} a \qty{60}{\degree}. O pequeno ângulo limite do diamante
aprisiona a luz que entra nele até ela sair pelas facetas de cima ---
o brilho.
\end{solution}

\begin{solution}{exo:b2:wave-interfaces:5}
(a) $\sqrt{1.52} = 1.23$; $e = 550/4 \times 1.38 = \qty{100}{nm}$. (b) $r_1 = -0.160$, $r_2 =
-0.048$, $\varphi = \pi$: $r = (r_1 - r_2)/(1 - r_1r_2) = -0.112$, $R = 1.3\%$ (contra
$4.3\%$). (c) A \qty{400}{nm}: $\varphi = 1.375\pi$, $R \approx 2.2\%$; a \qty{700}{nm}: $\varphi =
0.79\pi$, $R \approx 1.6\%$: o azul e o vermelho são mais refletidos que o verde ---
o brilho arroxeado de uma lente revestida. (d) Várias camadas achatam $R$
ao longo do visível.
\end{solution}

\begin{solution}{exo:b2:wave-interfaces:6}
(a) $T = 4 \times 30 \times 1.6/31.6^2 = 19\%$; o resto reverbera, e o pulso
ressoa. (b) $Z = \sqrt{Z_1Z_2} = \qty{6.9e6}{kg.m^{-2}.s^{-1}}$, $e = v/4f = \qty{0.15}{mm}$.
(c) $T_1 = 4 \times 30 \times 10^6 \times 410/(30 \times 10^6)^2 = 5.5 \times 10^{-5}$, $T_2 = 4 \times 1.6 \times
10^6 \times 410/(1.6 \times 10^6)^2 = 1.0 \times 10^{-3}$: $5.6 \times 10^{-8}$. (d) $4 \times 1.5 \times 1.6/
3.1^2 = 0.999$.
\end{solution}

\begin{solution}{exo:b2:wave-interfaces:7}
(a) $\omega_p/\omega = 3$: $\underline n = -\iu\sqrt8 = -2.83\iu$; $r = (1 + 2.83\iu)/(1 - 2.83
\iu)$, $|r| = 1$, fase $2\arctan2.83 = \qty{141}{\degree}$. (b) $c/\sqrt{\omega_p^2 - \omega^2}
= \qty{5.6}{m}$. (c) $E \propto \cos(kx + \varphi/2)$: primeiro nó em $kx = -(\pi + \varphi)/2$,
$x \approx \qty{-45}{m}$ abaixo da camada. (d) O índice cai com a altura, de modo que
Descartes afasta o raio da vertical; ele dá meia-volta onde
$n(z) = \sin\theta_0$, ou seja, $\omega_p^2 = \omega^2\cos^2\theta_0$, abaixo do nível em que $\omega = \omega_p$;
uma frequência cujo $\omega\cos\theta_0$ excede o $\omega_p$ máximo da camada
escapa.
\end{solution}

\begin{solution}{exo:b2:wave-interfaces:8}
$\sqrt{8\varepsilon_0\omega/\gamma}$: $2.7 \times 10^{-4}$ ($R = 0.9997$), $1.5 \times 10^{-2}$ ($0.985$),
$0.06$ ($0.94$) --- mas a \qty{5e14}{Hz} $\omega\tau = 80$: o cobre é um \omterm{def:b2:waves-in-media:plasma}{plasma},
e não um condutor ôhmico, e reflete cerca de $60\%$, com cor. \qty{15}{W}
absorvidos do feixe de CO$_2$: algumas centenas de kelvin por minuto num
espelho pequeno, daí a água.
\end{solution}

\begin{solution}{exo:b2:wave-interfaces:9}
$r = 0.2$, $-1/3$, $-1$, $+1$, $0$; \omterm{prop:b2:wave-interfaces:coax}{razão de onda estacionária} $1.5$, $2$, $\infty$, $\infty$, $1$; refletidos
$4\%$, $11\%$, $100\%$, $100\%$, $0$. \qty{96}{W} chegam à antena e \qty{4}{W}
voltam ao transmissor (e são em parte rerrefletidos). \qty{1.2}{\micro s}:
\qty{120}{m}; eco direito: um circuito aberto --- um cabo cortado.
\end{solution}

\begin{solution}{exo:b2:wave-interfaces:10}
(a) $\vect E = E\vect e_z$; $\vect k\wedge\vect e_z = (k_y, -k_x, 0)$, de modo que a componente
tangencial $B_y = -(n/c)\cos\theta\,E$ para as ondas incidente e transmitida e
$+(n_1/c)\cos\theta_1E_r$ para a refletida. (b) $1 + r = t$ e $n_1\cos\theta_1
(1 - r) = n_2\cos\theta_2t$: o $r_s$ enunciado. (c) $\theta = 0$: $(n_1 - n_2)/(n_1 +
n_2)$; $\theta_1 \to 90^\circ$: $r_s \to -1$. (d) $\cos\theta_2 = 0.882$: $r_s = -0.30$, $R_s =
9.2\%$; $R_p = 0.85\%$; contra $4.3\%$ na incidência normal.
\end{solution}

\begin{solution}{exo:b2:wave-interfaces:11}
(a) \qty{270}{nm}. (b) $d = 0.35\ell = \qty{95}{nm}$; $2.3\ell = \qty{620}{nm}$. (c) As
cristas tocam o vidro e deixam a luz sair para a pele: cristas
escuras e vales claros no detector. (d) A \omterm{def:b2:dispersion-wave-packets:complex}{onda evanescente} é a
função de onda que decai exponencialmente na região proibida; um segundo
meio dentro do comprimento de decaimento a recolhe --- tunelamento óptico.
\end{solution}

\begin{solution}{exo:b2:wave-interfaces:12}
(a) Nós de $\vect E$ nos dois espelhos: $\sin kL = 0$, $L = m\lambda/2$, $f_m = mc/2L$.
(b) $1$, $2$, \qty{3}{GHz}: nenhum em \qty{2.45}{GHz} --- os modos tridimensionais
de um forno real são muito mais densos. (c) $j_s = H_0 = E_0/\mu_0c = \qty{53}{A/m}$
em cada um, e a pressão $\mu_0j_s^2/2 = \qty{1.8}{mPa}$ afasta os espelhos
um do outro. (d) $\delta = \qty{1.3}{\micro m}$, $R_s = \qty{12.5}{m\ohm}$; $\tfrac12R_sH_0^2 =
\qty{18}{W/m^2}$ por espelho; energia média armazenada $\tfrac14\varepsilon_0E_0^2 \times L = \qty{1.3e-7}{J/m^2}$;
$Q = 2\pi f \times 1.3 \times 10^{-7}/35 \approx 6 \times 10^4$.
\end{solution}

\begin{solution}{pb:b2:wave-interfaces:1}
\textbf{1.} $R = 4.3\%$; $0.957^{16} = 0.50$; metade da luz se perde.

\textbf{2.} $n = 1.23$, $e = \qty{112}{nm}$; nenhum sólido durável tem índice
tão baixo (o MgF$_2$, $1.38$, é o menor): residual de $1.3\%$.

\textbf{3.} $0.987^{16} = 0.81$; $0.997^{16} = 0.95$. Os reflexos parasitas embaçam a
imagem: véu de luz, contraste baixo.

\textbf{4.} $\varphi = \pi \times 550/440 = 1.25\pi$: $R \approx 0.0256 + 0.0025 + 2 \times 0.008 \times
\cos1.25\pi = 1.7\%$ --- mais que a \qty{550}{nm}, como na ponta vermelha: um
reflexo magenta.

\textbf{5.} Cada interface reflete $|r| = 0.97/3.73 = 0.26$; o espaçamento de
quarto de onda põe as reflexões sucessivas em fase, de modo que as amplitudes se somam
em vez de se cancelar; a refletância da pilha vale $1 - 4(n_L/n_H)^{2N}
\approx 1 - 10^{-9}$ para $N = 20$.

\textbf{6.} A condição de estarem em fase só vale perto do comprimento de onda
de projeto (uma faixa de uns $\pm15\%$); o resto passa: um espelho verde
parece magenta por transmissão.

\textbf{7.} Para a onda transmitida: a interferência redistribui, com
$R + T = 1$.

\textbf{8.} $19\%$; com $Z_m = \sqrt{Z_1Z_2} = \qty{6.9e6}{kg.m^{-2}.s^{-1}}$: $100\%$ a
\qty{5}{MHz}.

\textbf{9.} $r_1 = 0.63$, $r_2 = 0.62$; a \qty{3}{MHz} a fase de ida e volta vale
$0.6\pi$, e a \qty{7}{MHz}, $1.4\pi$: $R \approx 0.39 + 0.39 - 0.24 = 0.54$, $T \approx 0.46$
nas duas: o casamento só vale numa faixa em torno de \qty{5}{MHz}, ao passo que um
pulso curto precisa de faixa larga --- um compromisso.

\textbf{10.} $t_1 = 2 \times 410/30 \times 10^6 = 2.7 \times 10^{-5}$, $t_2 = 2.0$: $t = 5.5 \times 10^{-5}$,
$T = t^2Z_1/Z_3 = 5.6 \times 10^{-8}$: \qty{-72}{dB}.

\textbf{11.} $t = 0.095 \times 1.03 = 0.098$, $T = 0.098^2 \times 30/1.6 = 0.18$: o gel
apenas elimina o ar.

\textbf{12.} $0.9^N = 0.01$: $N = 44$ idas e voltas --- um repique longo; o
apoio absorve a onda que volta e a mata.

\textbf{13.} $R = 0.40$ no osso; $2 \times 5 \times 0.5 \times 5 = \qty{25}{dB}$ de
\omterm{def:b2:dispersion-wave-packets:complex}{atenuação}: $0.40 \times 3.2 \times 10^{-3} = 1.3 \times 10^{-3}$ (\qty{-29}{dB}); o osso
absorve o que não reflete: nada volta de trás dele.

\textbf{14.} $R = 1$, $r = -1$.

\textbf{15.} Curto em $x = 0$: $v = -2\iu v_i\sin kx$, $i = (2v_i/Z_0)\cos kx$;
em $x = -d$: $v/i = \iu Z_0\tan kd$; infinita em $d = \lambda/4$.

\textbf{16.} $R_s \parallel \infty = R_s$: $r = (R_s - Z_0)/(R_s + Z_0) = 0$ para
\qty{377}{\ohm}; $d = \qty{7.5}{mm}$.

\textbf{17.} A lâmina vê o campo incidente inteiro: $|v|^2/2R_s =
E_0^2/2Z_0$, a intensidade incidente --- tudo absorvido.

\textbf{18.} A \qty{8}{GHz}: $kd = 72^\circ$, $Z_{\text{ent}} = \iu\qty{1160}{\ohm}$, $R_s \parallel
Z_{\text{ent}} = (340 + 111\iu)\,\Omega$, $R = |r|^2 = 2.6\%$ (\qty{-16}{dB}); a \qty{12}{GHz}:
$kd = 108^\circ$, mesmo módulo: \qty{-16}{dB}. Útil em cerca de $\pm20\%$.

\textbf{19.} A incidência oblíqua encurta o quarto de onda efetivo
($d\cos\theta$) e muda a impedância de onda com o ângulo: o
casamento se perde.

\textbf{20.} Uma tela funciona numa frequência e num ângulo; camadas
graduadas com perdas e superfícies moldadas absorvem e desviam numa banda larga.

\textbf{21.} $\underline n = -\iu\sqrt{(\omega_p/\omega)^2 - 1} = -3.6\iu$; $|r| = 1$, fase
$2\arctan3.6 = \qty{149}{\degree}$; profundidade $c/3.6\omega = \qty{22}{nm}$. A onda nunca
entra no metal para ser dividida em duas reflexões comparáveis: não há
nada contra o que uma \omterm{prop:b2:wave-interfaces:optics}{camada de quarto de onda} possa cancelar uma reflexão
de módulo um.

\textbf{22.} \qty{30}{W}; $30/(0.01 \times 235) = \qty{13}{K/s}$.

\textbf{23.} Nós a cada $\lambda/2 = \qty{250}{nm}$; $E \propto \cos(kx + \varphi/2)$ com $\varphi =
149^\circ$: primeiro nó em $x = -(90^\circ + 74.5^\circ)\lambda/360^\circ = \qty{-230}{nm}$; um
grão ali fica no escuro.

\textbf{24.} $I = \qty{1e7}{W/m^2}$, $E_0 = \sqrt{2I/\varepsilon_0c} = \qty{87}{kV/m}$, $j_s =
2E_0/\mu_0c = \qty{460}{A/m}$.

\textbf{25.} Metal: \omterm{prop:b2:maxwell-equations:boundary}{correntes superficiais} numa \omterm{def:b2:dispersion-wave-packets:complex}{profundidade de penetração}, $r \approx -1$, limitado
pela perda Joule dessas correntes. \omterm{def:b2:waves-in-media:plasma}{Plasma}: penetração evanescente,
$|r| = 1$ com fase dependente da frequência, limitado pelas colisões.
Pilha dielétrica: muitas reflexões fracas somadas em fase, $r \to 1$ numa
banda, limitado pelo número de camadas e pela largura de banda.
\end{solution}
